\documentclass[11pt,a4paper]{article}

\usepackage[T1]{fontenc}
\usepackage[margin=1in]{geometry}
\usepackage[authoryear,longnamesfirst]{natbib}
\usepackage{amsmath,amssymb,amsthm,mathtools}
\usepackage{enumitem}
\usepackage{booktabs}
\usepackage{array}
\usepackage{tabularx}
\usepackage{xurl}
\usepackage{microtype}
\usepackage[hidelinks]{hyperref}
\hypersetup{
  pdftitle={A Reduction-Theoretic Formalization of Groebner Bases and an Abstract Buchberger Procedure in Lean 4},
  pdfauthor={Sanghyeok Park},
  pdfkeywords={Groebner bases, Buchberger criterion, polynomial reduction, abstract rewriting, Lean 4, Mathlib}
}
% Compact Lean 4 listing style for the manuscript.
\usepackage{xcolor}
\usepackage{listings}

\definecolor{LeanKeyword}{RGB}{0,70,140}
\definecolor{LeanTactic}{RGB}{120,40,120}
\definecolor{LeanComment}{RGB}{0,105,55}
\definecolor{LeanNumber}{RGB}{105,105,105}
\definecolor{LeanSymbol}{RGB}{0,0,130}

\lstdefinelanguage{Lean4}{
  sensitive=true,
  morekeywords=[1]{
    import,prelude,namespace,section,end,open,export,include,omit,
    variable,universe,noncomputable,private,protected,public,local,
    def,abbrev,theorem,lemma,example,structure,class,instance,inductive,
    where,deriving,mutual,partial,termination_by,decreasing_by
  },
  morekeywords=[2]{Sort,Type,Prop},
  morekeywords=[3]{
    by,fun,match,with,if,then,else,let,in,have,show,suffices,
    exact,apply,refine,intro,rintro,rcases,obtain,cases,induction,
    constructor,simp,simpa,rw,by_cases,change,clear
  },
  morecomment=[l]{--},
  morecomment=[s]{/-}{-/},
  morestring=[b]",
  literate=
    {α}{{$\alpha$}}1
    {σ}{{$\sigma$}}1
    {ℕ}{{$\mathbb{N}$}}1
    {∀}{{\color{LeanSymbol}$\forall$}}1
    {∃}{{\color{LeanSymbol}$\exists$}}1
    {→}{{\color{LeanSymbol}$\to$}}1
    {↔}{{\color{LeanSymbol}$\leftrightarrow$}}1
    {⟶}{{\color{LeanSymbol}$\longrightarrow$}}1
    {⟷}{{\color{LeanSymbol}$\longleftrightarrow$}}1
    {∈}{{\color{LeanSymbol}$\in$}}1
    {∉}{{\color{LeanSymbol}$\notin$}}1
    {⊆}{{\color{LeanSymbol}$\subseteq$}}1
    {≠}{{\color{LeanSymbol}$\neq$}}1
    {¬}{{\color{LeanSymbol}$\neg$}}1
    {∧}{{\color{LeanSymbol}$\wedge$}}1
    {∨}{{\color{LeanSymbol}$\vee$}}1
    {₀}{{$_0$}}1
    {₁}{{$_1$}}1
    {₂}{{$_2$}}1
    {:=}{{\color{LeanSymbol}:=}}2
}

\lstdefinestyle{lean}{
  language=Lean4,
  basicstyle=\ttfamily\footnotesize,
  keywordstyle=[1]\color{LeanKeyword}\bfseries,
  keywordstyle=[2]\color{LeanKeyword},
  keywordstyle=[3]\color{LeanTactic},
  commentstyle=\color{LeanComment}\itshape,
  numbers=left,
  numberstyle=\scriptsize\color{LeanNumber},
  numbersep=7pt,
  frame=single,
  framesep=4pt,
  columns=fullflexible,
  keepspaces=true,
  showstringspaces=false,
  breaklines=true,
  breakatwhitespace=false,
  captionpos=b
}


\newcommand{\lean}[1]{\nolinkurl{#1}}
\newcommand{\doclinkmark}{\textsuperscript{\ensuremath{\nearrow}}}
\newcommand{\leandoc}[3]{%
  \href{https://sanghyeok0.github.io/Buchberger/docs/#1.html\##2}{\lean{#3}}}
\newcommand{\K}{\mathbb{K}}
\newcommand{\N}{\mathbb{N}}
\newcommand{\supp}{\operatorname{supp}}
\newcommand{\LT}{\operatorname{LT}}
\newcommand{\LM}{\operatorname{LM}}
\newcommand{\LC}{\operatorname{LC}}
\newcommand{\spoly}{\operatorname{S}}
\newcommand{\ideal}[1]{\langle #1 \rangle}
\newcommand{\red}[1]{\xrightarrow[#1]{}}
\newcommand{\redstar}[1]{\xrightarrow[#1]{*}}
\newcommand{\eqvred}[1]{\xleftrightarrow[#1]{*}}
\newcommand{\joinred}[1]{\mathrel{\downarrow}_{#1}}
\newcommand{\Rpoly}{R[X_\sigma]}

\theoremstyle{plain}
\newtheorem{theorem}{Theorem}[section]
\newtheorem{proposition}[theorem]{Proposition}
\newtheorem{lemma}[theorem]{Lemma}
\newtheorem{corollary}[theorem]{Corollary}
\theoremstyle{definition}
\newtheorem{definition}[theorem]{Definition}
\newtheorem{remark}[theorem]{Remark}
\newtheorem{example}[theorem]{Example}

\title{A Reduction-Theoretic Formalization of Gr\"obner Bases and an Abstract Buchberger Procedure in Lean 4}
\author{Sanghyeok Park\\
Seoul National University, Seoul, Republic of Korea}
\date{}

\begin{document}
\maketitle

\begin{abstract}
We present a Lean 4 formalization of polynomial reduction, reduction-theoretic
characterizations of Gr\"obner bases, Buchberger's criterion, and an
abstract Buchberger procedure.  Reduction modulo a set $G$ is represented as a binary
relation on Mathlib's multivariate polynomials, providing a common interface
for finite reduction, normal forms, joinability, confluence, and the
Church--Rosser property.  We prove termination by a well-founded
colexicographic order on finite supports.  Assuming that every nonzero element
of $G$ has unit leading coefficient, we show that two polynomials are reduction
equivalent modulo $G$ exactly when their difference belongs to the ideal
$\ideal{G}$.  This gives equivalent characterizations of the Gr\"obner
condition by reduction to zero for all elements of $\ideal{G}$,
Church--Rosser, local confluence, confluence, and uniqueness of normal forms.
Assuming additionally that the coefficient ring has no zero divisors, an
analysis of pairs of one-step reductions from the same polynomial proves
Buchberger's criterion.  Finally, we define a nondeterministic Buchberger step
relation that adjoins nonzero normal forms of S-polynomials.  Every Buchberger
step preserves the generated ideal.  After the unit-leading-coefficient
invariant is incorporated into the state space, every state step strictly
enlarges the ideal generated by the current leading terms.  If the polynomial
ring is Noetherian and the coefficient ring is a field, every initial set
reaches a terminal set that is a Gr\"obner basis of the ideal it generates.  In
the finite-variable field specialization, finite initial sets reach finite
terminal Gr\"obner bases.  The construction is a verified specification rather
than an executable algorithm with a fixed pair-selection strategy.
\end{abstract}

\noindent\textbf{Keywords:} Gr\"obner bases, Buchberger criterion, polynomial reduction,
abstract rewriting, formal verification, Lean 4, Mathlib

\medskip
\section{Introduction}\label{sec:introduction}

Fix a monomial order $m$ on a multivariate polynomial ring.  Throughout
this article, all leading data, S-polynomials, and polynomial reductions are
taken with respect to $m$, and this dependence is suppressed from the
mathematical notation.  Let $I$ be an ideal and let $G\subseteq I$.  The set
$G$ is a Gr\"obner basis of $I$ when
\[
  \ideal{\LT(g)\mid g\in G}
  =
  \ideal{\LT(f)\mid f\in I}.
\]
The ideal on the right is the leading-term ideal of $I$.
This algebraic equality is connected with algorithmic properties of polynomial
reduction, including reduction to normal form, reduction of S-polynomials, and
the abstract Buchberger procedure
\citep{Buchberger2006,BeckerWeispfenning1993,CoxLittleOShea2025}.
A formal development must state the reduction relation and prove each
connection with this equality of ideals explicitly.

This article develops that reduction-theoretic account in Lean 4
\citep{deMouraUllrich2021}.  The development builds on Mathlib's
infrastructure for multivariate polynomials, monomial orders, leading terms,
S-polynomials, and ideals \citep{Mathlib2020}.  It also proves bridge lemmas translating finite relational reduction
sequences into finite linear-combination certificates satisfying the same
product-degree bound that appears in Mathlib's division theorem.  Its basic reduction relation allows one step to cancel any
reducible monomial in a polynomial.
Because the reduction operation is represented as a binary relation, the
standard definitions of finite reduction, joinability, equivalence closure,
normal forms, local confluence, confluence, and Church--Rosser apply without
being tied to a particular division strategy.

The formalization uses two predicates in the namespace
\lean{MonomialOrder}.  The predicate
\leandoc{Mathlib/RingTheory/MvPolynomial/GroebnerBasisCriterion}{MonomialOrder.IsGroebner}{IsGroebner}
compares the leading
terms of $G$ with the leading-term ideal of $\ideal{G}$.  The predicate
\leandoc{Mathlib/RingTheory/MvPolynomial/BuchbergerCriterion}{MonomialOrder.IsGroebnerBasis}{IsGroebnerBasis}
states that $G\subseteq I$ and compares the leading
terms of $G$ with the leading-term ideal of the specified ideal $I$.  The two
predicates are equivalent when $I=\ideal{G}$.  Throughout the article,
statements about \lean{IsGroebner} are called \emph{self-relative}, whereas
statements about \lean{IsGroebnerBasis} are called \emph{ideal-relative}.

The main contributions are as follows.
\begin{enumerate}[label=(\roman*),leftmargin=2.2em]
  \item We extend Mathlib's relation-closure API with explicit definitions of
    confluence and Church--Rosser, and we formalize normal forms, local
    confluence, Newman's lemma, and uniqueness of normal forms.
  \item We define arbitrary-term polynomial reduction and prove that its
    one-step relation admits no infinite forward reduction chain.  The proof
    uses a well-founded colexicographic order on the finite supports of
    polynomials.
  \item Assuming that every nonzero $g\in G$ has unit leading coefficient,
    we prove the exact translation
    \[
      f\eqvred{G}g
      \quad\Longleftrightarrow\quad
      f-g\in\ideal{G}.
    \]
  \item Under the same unit-leading-coefficient hypothesis, we prove that
    \[
      \ideal{\LT(f)\mid f\in\ideal{G}}
      = \ideal{\LT(g)\mid g\in G}
    \]
    is equivalent to reduction to zero for every element of $\ideal{G}$,
    Church--Rosser, local confluence, confluence, and uniqueness of normal
    forms.
  \item Assuming additionally that $R$ has no zero divisors, we prove
    Buchberger's criterion by showing that every local peak formed by two
    one-step reductions is joinable: the same-exponent case is controlled by
    an S-polynomial, while the distinct-exponent case is resolved by a direct
    reduction of the difference to zero.
  \item We formalize an abstract Buchberger procedure as a nondeterministic
    step relation.  Every step preserves the generated ideal.  After the
    unit-leading-coefficient invariant is incorporated into the state space,
    each state step strictly increases the current leading-term ideal; hence
    the state relation has no infinite forward chain over a Noetherian
    polynomial ring.  If $R$ has no zero divisors and every nonzero polynomial
    has unit leading coefficient, a terminal state is a Gr\"obner basis of
    the ideal generated by the initial polynomial set.  One step and every
    finite chain of steps preserve finiteness; consequently, in the
    finite-variable field specialization, finite input yields a finite terminal
    Gr\"obner basis.
\end{enumerate}

The phrase \emph{abstract Buchberger procedure} refers to the
nondeterministic relation defined in Section~\ref{sec:buchberger-procedure}.
It is a predicate on pairs of states: it specifies when an enlarged polynomial
set is an admissible successor of the current set, rather than defining a
function that computes one particular successor.  In particular, the
formalization does not maintain a finite queue of unprocessed critical pairs,
choose a pair by a deterministic selection function, or compute a normal form
by an executable procedure.  Under the hypotheses stated above, the state
relation admits no infinite forward chain, and every reachable terminal state
is correct; however, the formalization does not provide a function that can be
evaluated on a concrete input to return a Gr\"obner basis.
Sections~\ref{sec:related-work} and \ref{sec:future-work} compare this
relational specification with executable verified algorithms and with tactics
that obtain candidates from external computer algebra systems.

\subsection{Source organization}

The canonical Lean source is maintained in the
\lean{buchberger-formalization} branch of the author's \lean{mathlib4} fork,
since the development extends Mathlib modules directly and is intended for
upstream integration.

A separate \lean{Buchberger} repository serves as the companion artifact and
documentation site.  It does not duplicate the Lean source.  Instead, its root
module \lean{Buchberger.lean} imports the six modules from the pinned
\lean{mathlib4} snapshot.  Its Lake configuration fixes the source dependency
and Lean toolchain, while GitHub Actions verify the build and generate
\lean{doc-gen4} API documentation with declaration-level links back to the
source.

The development consists of five new files, together with additional
declarations in Mathlib's existing relation file.
Figure~\ref{fig:dependencies} shows the import dependency used in the
formalization.

\begin{figure}[t]
\centering
\setlength{\tabcolsep}{2pt}
\begin{tabular}{c}
\fbox{\strut\lean{Relation.lean}}\\[-1pt]
$\downarrow$\\[-1pt]
\fbox{\strut\lean{NormalForm.lean}}\\[-1pt]
$\downarrow$\\[-1pt]
\fbox{\strut\lean{PolynomialReductions.lean}}\\[-1pt]
$\downarrow$\\[-1pt]
\fbox{\strut\lean{GroebnerBasisCriterion.lean}}\\[-1pt]
$\downarrow$\\[-1pt]
\fbox{\strut\lean{BuchbergerCriterion.lean}}\\[-1pt]
$\downarrow$\\[-1pt]
\fbox{\strut\lean{BuchbergerAlgorithm.lean}}
\end{tabular}
\caption{Import dependency of the formal development.  The box labelled
\lean{Relation.lean} refers only to the new confluence and Church--Rosser
declarations added to that existing Mathlib file.  The other five files
constitute the reduction-theoretic Gr\"obner basis development described in
this article.}
\label{fig:dependencies}
\end{figure}

Table~\ref{tab:inventory} lists the declarations corresponding to the
principal theorems below.  In this table and
Table~\ref{tab:groebner-characterizations-lean}, \doclinkmark marks links to
the generated API documentation.

\begin{table}[t]
\caption{Selected mathematical results and their Lean declarations.}
\label{tab:inventory}
\centering
\scriptsize
\begin{tabularx}{\textwidth}{@{}>{\raggedright\arraybackslash}p{0.20\textwidth}>{\raggedright\arraybackslash}p{0.54\textwidth}>{\raggedright\arraybackslash}X@{}}
\toprule
Mathematical result & Lean declaration \doclinkmark & Source module \\
\midrule
Newman's lemma & \leandoc{Mathlib/Logic/Relation/NormalForm}{Relation.confluent_of_wellFounded_flip_of_locallyConfluent}{Relation.confluent_of_wellFounded_flip_of_locallyConfluent} & \lean{NormalForm} \\
Reduction termination & \leandoc{Mathlib/RingTheory/MvPolynomial/PolynomialReductions}{MonomialOrder.ReducesToSet.wellFounded}{MonomialOrder.ReducesToSet.wellFounded} & \lean{PolynomialReductions} \\
Finite reduction certificate with degree bound & \leandoc{Mathlib/RingTheory/MvPolynomial/PolynomialReductions}{MonomialOrder.exists_linearCombination_of_reflTransGen_with_degree_bound}{MonomialOrder.exists_linearCombination_of_reflTransGen_with_degree_bound} & \lean{PolynomialReductions} \\
Reduction equivalence and ideal congruence & \leandoc{Mathlib/RingTheory/MvPolynomial/PolynomialReductions}{MonomialOrder.eqvGen_iff_sub_mem_span}{MonomialOrder.eqvGen_iff_sub_mem_span} & \lean{PolynomialReductions} \\
Gröbner/confluence equivalence & \leandoc{Mathlib/RingTheory/MvPolynomial/GroebnerBasisCriterion}{MonomialOrder.IsGroebner.iff_confluent}{MonomialOrder.IsGroebner.iff_confluent} & \lean{GroebnerBasisCriterion} \\
Buchberger criterion & \leandoc{Mathlib/RingTheory/MvPolynomial/BuchbergerCriterion}{MonomialOrder.IsGroebner.iff_sPolynomial_reflTransGen_zero}{MonomialOrder.IsGroebner.iff_sPolynomial_reflTransGen_zero} & \lean{BuchbergerCriterion} \\
Termination of Buchberger steps & \leandoc{Mathlib/RingTheory/MvPolynomial/BuchbergerAlgorithm}{MonomialOrder.BuchbergerState.wellFounded_flip}{MonomialOrder.BuchbergerState.wellFounded_flip} & \lean{BuchbergerAlgorithm} \\
Finite terminal Gr\"obner basis & \leandoc{Mathlib/RingTheory/MvPolynomial/BuchbergerAlgorithm}{MonomialOrder.BuchbergerState.exists_finite_isGroebnerBasis_terminal_field_of_finite_variables}{MonomialOrder.BuchbergerState.exists_finite_isGroebnerBasis_terminal_field_of_finite_variables} & \lean{BuchbergerAlgorithm} \\
\bottomrule
\end{tabularx}
\end{table}

\section{Preliminaries and abstract reduction systems}\label{sec:setting}

\subsection{Polynomial setting}

Let $R$ be a commutative ring, let the variables be indexed by $\sigma$, and write $\Rpoly$ for the multivariate polynomial ring.  
Lean represents this ring by \texttt{MvPolynomial }\(\sigma\)\texttt{ R}.  
A monomial exponent is a finitely supported function $\alpha:\sigma\to_0\N$.  
We continue to work with the fixed monomial order $m$.  For every nonzero $f\in\Rpoly$, write $\LM(f)$, $\LC(f)$, and $\LT(f)$
for its leading monomial, leading coefficient, and leading term.  Thus
$\LT(f)=\LC(f)\LM(f)$.  Mathlib represents the exponent of $\LM(f)$ by
\lean{MonomialOrder.degree}; its leading coefficient and leading term are
\lean{MonomialOrder.leadingCoeff} and \lean{MonomialOrder.leadingTerm}.

\begin{remark}[Mathlib's convention for the zero polynomial]
\label{rem:zero-polynomial}
Mathlib makes its leading-data operations total.  In particular,
\[
  \texttt{m.degree 0}=0,\qquad \LC(0)=0,\qquad \LT(0)=0.
\]
For a nonzero polynomial, these operations have their usual meanings.  The
value assigned by \lean{m.degree} at zero is only a totalization convention;
it does not assert that the zero polynomial has a leading monomial.  Since
$\LT(0)=0$, including the zero polynomial among the generators of a
leading-term ideal does not change that ideal.
\end{remark}

\begin{definition}[Leading-term ideal and Gr\"obner conditions]
\label{def:groebner-conditions}
For an ideal $I\subseteq\Rpoly$, its leading-term ideal with respect to $m$ is
\[
  \ideal{\LT(f)\mid f\in I}.
\]
For a set $G\subseteq\Rpoly$, write $\ideal{G}$ for
$\operatorname{Ideal.span}(G)$.
\begin{enumerate}[label=(\roman*),leftmargin=2em]
  \item The \emph{self-relative Gr\"obner condition} for $G$ is
  \[
    \ideal{\LT(f)\mid f\in\ideal{G}}
    =\ideal{\LT(g)\mid g\in G}.
  \]
  \item If $G\subseteq I$, then $G$ is a \emph{Gr\"obner basis of $I$} when
  \[
    \ideal{\LT(f)\mid f\in I}
    =\ideal{\LT(g)\mid g\in G}.
  \]
\end{enumerate}
These conditions are represented by \lean{IsGroebner} and
\lean{IsGroebnerBasis} in the namespace \lean{MonomialOrder}, respectively.
\end{definition}

\begin{definition}[S-polynomial]
\label{def:s-polynomial}
Let $f,g\in\Rpoly$ be nonzero, and let
\[
  M=\operatorname{lcm}\bigl(\LM(f),\LM(g)\bigr)
\]
be the least common multiple of their leading monomials.  The
\emph{S-polynomial} of $f$ and $g$ is
\begin{equation}\label{eq:spoly}
  \spoly(f,g)
  =
  \LC(g)\frac{M}{\LM(f)}f
  -
  \LC(f)\frac{M}{\LM(g)}g.
\end{equation}
Here the quotients are monomial quotients.  The leading coefficients are
cross-multiplied, so Formula~\eqref{eq:spoly} does not require coefficient
inverses.
\end{definition}

\begin{remark}[Mathlib's total definition of the S-polynomial]
\label{rem:spoly-total}
Definition~\ref{def:s-polynomial} is the mathematical description for
nonzero inputs.  Mathlib defines \lean{MonomialOrder.sPolynomial} for all
polynomials.  Rather than explicitly forming $M$, Mathlib represents the two
quotient exponents by the componentwise truncated differences
$\deg_m(g)-\deg_m(f)$ and $\deg_m(f)-\deg_m(g)$.  These are equal,
respectively, to
\[
  \bigl(\deg_m(f)\sqcup\deg_m(g)\bigr)-\deg_m(f)
  \quad\text{and}\quad
  \bigl(\deg_m(f)\sqcup\deg_m(g)\bigr)-\deg_m(g),
\]
and hence agree with the monomial quotients in
Formula~\eqref{eq:spoly}.  Together with the conventions in
Remark~\ref{rem:zero-polynomial}, the total definition also gives
\[
  \spoly(0,g)=0,\qquad
  \spoly(f,0)=0,\qquad
  \spoly(f,f)=0.
\]
\end{remark}

\subsection{Abstract reduction systems}\label{sec:abstract-reduction}

Let $r:\alpha\to\alpha\to\mathrm{Prop}$, and write $r^*$ for its
reflexive-transitive closure.

\begin{definition}[Normal forms and uniqueness]
\label{def:normal-forms}
An element $b$ is a \emph{normal form} if it has no outgoing $r$-step.  It is
a \emph{normal form of $a$} if $r^*(a,b)$ and $b$ is a normal form.  The
relation $r$ has \emph{unique normal forms} if any two normal forms reachable
from the same element are equal.
\end{definition}

\begin{definition}[Joinability and confluence properties]
\label{def:confluence-properties}
Elements $b$ and $c$ are \emph{joinable} if there is a $d$ such that
$r^*(b,d)$ and $r^*(c,d)$.
\begin{enumerate}[label=(\roman*),leftmargin=2em]
  \item The relation $r$ is \emph{locally confluent} if the two targets of
        every pair of one-step reductions from the same source are joinable.
  \item It is \emph{confluent} if the two targets of every pair of finite
        reductions from the same source are joinable.
  \item It has the \emph{Church--Rosser property} if any two elements related
        by the equivalence closure of $r$ are joinable.
\end{enumerate}
\end{definition}

Listing~\ref{lst:normal-forms} gives the corresponding Lean definitions of
normal forms, uniqueness of normal forms, and local confluence.

\begin{lstlisting}[style=lean,caption={Normal forms and local confluence for an abstract relation.},label={lst:normal-forms}]
def IsNormalForm (r : α → α → Prop) (a : α) : Prop :=
  ∀ b, ¬ r a b

def IsNormalFormOf (r : α → α → Prop) (a b : α) : Prop :=
  ReflTransGen r a b ∧ IsNormalForm r b

def UniqueNormalForms (r : α → α → Prop) : Prop :=
  ∀ {a b c : α},
    ReflTransGen r a b → ReflTransGen r a c →
    IsNormalForm r b → IsNormalForm r c → b = c

def LocallyConfluent (r : α → α → Prop) : Prop :=
  ∀ {a b c : α},
    r a b → r a c → Join (ReflTransGen r) b c
\end{lstlisting}

A relation $r$ is called \emph{terminating} if there is no infinite
sequence $a_0,a_1,a_2,\ldots$ satisfying
\[
  r(a_i,a_{i+1})
  \qquad\text{for every }i\geq0.
\]
Lean encodes this condition as well-foundedness of $\operatorname{flip}(r)$:
a forward $r$-sequence becomes a descending chain for the flipped relation.

\begin{theorem}[Existence of normal forms]\label{thm:normal-form-exists}
If $r$ is terminating, then every element has a normal form.
\end{theorem}

\begin{theorem}[Newman's lemma \citep{Newman1942}]\label{thm:newman}
If $r$ is terminating and locally confluent, then $r$ is confluent.
\end{theorem}

\begin{proposition}[Confluence, Church--Rosser, and normal forms]
\label{prop:abstract-equivalences}
For every relation $r$, confluence implies uniqueness of normal forms, and
confluence is equivalent to the Church--Rosser property.  If $r$ is
terminating, then uniqueness of normal forms also implies confluence.
\end{proposition}

Theorems~\ref{thm:normal-form-exists}--\ref{thm:newman} and
Proposition~\ref{prop:abstract-equivalences} are standard results about
abstract reduction systems; see \citet[Chapter~2]{BaaderNipkow1998}.
In particular, the Church--Rosser characterization and uniqueness of normal
forms are discussed in Section~2.1, while Newman's lemma appears as
Lemma~2.7.2. They form the abstract rewriting layer used below and do not
depend on polynomials.

\section{Polynomial reduction}\label{sec:polynomial-reduction}

\subsection{One-step reduction}

Continue to use the fixed monomial order $m$.

\begin{definition}[One-step polynomial reduction]
\label{def:one-step-reduction}
Let $p,f,g\in\Rpoly$.  We say that $f$ reduces to $g$ by $p$ if
there are monomials $x^t$ and $x^s$ and a coefficient $c\in R$ such that
\begin{enumerate}[label=(\roman*),leftmargin=2em]
  \item $p\neq0$ and $x^t$ occurs in $f$;
  \item $\LM(p)\mid x^t$; equivalently, $x^t=x^s\LM(p)$;
  \item $c\LC(p)$ equals the coefficient of $x^t$ in $f$; and
  \item $g=f-cx^s p$.
\end{enumerate}
Reduction modulo a set $G$ means reduction by some polynomial $p\in G$.
The Lean definition records the exponents $s$ and $t$ explicitly and writes
the coefficient of $x^t$ in $f$ as \lean{f.coeff t}.
\end{definition}

The exact Lean definition is shown in
Listing~\ref{lst:reduction-definition}.

\begin{lstlisting}[style=lean,caption={One-step polynomial reduction and its existential closures.},label={lst:reduction-definition}]
def ReducesToBy
    (p f g : MvPolynomial σ R) (t : σ →₀ ℕ) : Prop :=
  p ≠ 0 ∧
    t ∈ f.support ∧
      ∃ s : σ →₀ ℕ,
        s + m.degree p = t ∧
          ∃ c : R,
            c * m.leadingCoeff p = f.coeff t ∧
              g = f - monomial s c * p

def ReducesToPoly (p f g : MvPolynomial σ R) : Prop :=
  ∃ t : σ →₀ ℕ, m.ReducesToBy p f g t

def ReducesToSet
    (P : Set (MvPolynomial σ R))
    (f g : MvPolynomial σ R) : Prop :=
  ∃ p ∈ P, m.ReducesToPoly p f g
\end{lstlisting}

\begin{remark}[Single-reducer strong reduction]
\label{rem:strong-reduction}
The relation in Definition~\ref{def:one-step-reduction} is a single-reducer
strong reduction relation: one step subtracts a monomial multiple of one
polynomial $p\in G$.  It does not implement weak reduction over coefficient
rings, where leading coefficients of several reducers may be combined to
cancel a single term.  For example, in $\mathbb{Z}[x]$, let
\[
  G=\{2x,3x\},\qquad f=x.
\]
Under the present relation, $x$ is normal modulo $G$, because neither $2$ nor
$3$ divides its coefficient $1$.  In a weak reduction framework, however, the
identity
\[
  x=(-1)(2x)+(3x)
\]
allows the term to be cancelled by combining the two reducers.
\end{remark}

\begin{definition}[Reduction closures]
\label{def:reduction-closures}
For a fixed set $G\subseteq\Rpoly$:
\begin{enumerate}[label=(\roman*),leftmargin=2em]
  \item $f\red{G}g$ denotes one-step reduction by some $p\in G$;
  \item $f\redstar{G}g$ denotes the reflexive-transitive closure of
        one-step reduction;
  \item $f\eqvred{G}g$ denotes the smallest equivalence relation containing
        one-step reduction; and
  \item $f\joinred{G}g$ means that there exists $h$ with
        $f\redstar{G}h$ and $g\redstar{G}h$.
\end{enumerate}
\end{definition}

The relation $\red{G}$ permits cancellation of any reducible monomial.
Mathlib's division procedure, by contrast, selects terms according to a
specific algorithm.  The arbitrary-term relation is used for confluence and
Church--Rosser statements; leading-term reduction remains available as the
separate predicate \lean{LTermReducesToPoly}.

\subsection{Reducibility over rings}

\begin{proposition}[Reducibility criterion over a commutative ring]
\label{prop:reducibility-ring}
For $p,f\in\Rpoly$, the polynomial $f$ is reducible by $p$ if and only if
\[
  p\neq0
  \quad\text{and}\quad
  \exists t\in\supp(f),\qquad
    \LM(p)\mid x^t
    \quad\text{and}\quad
    \LC(p)\mid\operatorname{coeff}_t(f).
\]
Since $\LM(p)=x^{\deg_m(p)}$, the monomial-divisibility condition is
equivalently the componentwise inequality $\deg_m(p)\leq t$ on exponent
vectors.  Thus reducibility over a general coefficient ring consists of both
monomial divisibility and coefficient divisibility.
\end{proposition}

Proposition~\ref{prop:reducibility-ring} is the exact content of
\lean{MonomialOrder.ReducesToPoly.reducible_iff_exists_degree_le_and_leadingCoeff_dvd}.
In particular, a unit leading coefficient is not required to define a
reduction step.  It is a sufficient condition that makes the coefficient
divisibility requirement automatic and reduces the criterion to monomial
divisibility alone.

\begin{remark}[The unit-leading-coefficient hypothesis]
\label{rem:unit-condition}
Several later translation and confluence theorems assume
\begin{equation}\label{eq:unit-condition}
  \forall g\in G,\qquad
  g\neq0\;\Longrightarrow\;\LC(g)\text{ is a unit}.
\end{equation}
The Lean statements encode the same condition as
\(
  \operatorname{IsUnit}(\LC(g))\lor g=0
\), which avoids imposing a unit condition on \(\LC(0)=0\).  Under this
hypothesis, whenever $\LM(g)$ divides a selected monomial, the required
coefficient multiple exists.  Over a field, Condition~\eqref{eq:unit-condition}
is automatic.  Over a general ring, it is an algebraic hypothesis used by
specific theorems rather than by the definition of reduction itself.
\end{remark}

\subsection{Termination by support colex order}

Total degree is not a termination measure for arbitrary-term reduction:
cancelling one monomial can introduce several monomials of the same total
degree.

\begin{definition}[Support colexicographic order]
\label{def:support-colex}
Identify the support of a polynomial with the finite set of monomials that
occur with nonzero coefficient.  If
$\supp(f)\mathbin{\triangle}\supp(g)\neq\varnothing$, let
\[
  M=\max\nolimits_m\bigl(\supp(f)\mathbin{\triangle}\supp(g)\bigr)
\]
be the largest monomial in the symmetric difference with respect to the fixed
monomial order.  We declare
\[
  f<_{\mathrm{supp}}g
\]
when $M\in\supp(g)$.  If the supports are equal, their symmetric difference is
empty and neither polynomial is smaller than the other.  Thus the polynomial
containing the largest monomial at which the two supports differ is the larger
one.
\end{definition}

In Lean, monomials are stored as exponent vectors.  The implementation
transports these exponents to the linearly ordered synonym attached to $m$
and then applies the finite-set colexicographic order; this is the role of
\lean{m.toSyn} and \lean{toColex} in Listing~\ref{lst:termination}.

\begin{lstlisting}[style=lean,caption={The support-colex measure and termination of polynomial reduction.},label={lst:termination}]
def supportColexLt (f g : MvPolynomial σ R) : Prop :=
  letI : LinearOrder m.syn := m.linearOrderSyn
  toColex (f.support.image m.toSyn) <
    toColex (g.support.image m.toSyn)

theorem supportColexLt_wellFounded :
    WellFounded
      (m.supportColexLt :
        MvPolynomial σ R → MvPolynomial σ R → Prop)

theorem ReducesToSet.wellFounded
    (P : Set (MvPolynomial σ R)) :
    WellFounded (flip (m.ReducesToSet P)) :=
  m.supportColexLt_wellFounded.mono
    (by
      intro f g h
      exact ReducesToSet.supportColexLt m P h)
\end{lstlisting}

Suppose that $f\red{G}g$ cancels the monomial $x^t$.  Then
$x^t\in\supp(f)\setminus\supp(g)$, while every monomial that occurs in $g$
but not in $f$ is strictly smaller than $x^t$.  Hence $x^t$ is the largest
monomial at which the supports differ, and it belongs to $f$.  Therefore
$g<_{\mathrm{supp}}f$.

\begin{theorem}[Termination of polynomial reduction]\label{thm:reduction-termination}
For every set $G\subseteq\Rpoly$, there is no infinite forward reduction
chain
\[
  f_0\red{G}f_1\red{G}f_2\red{G}\cdots.
\]
\end{theorem}

Listing~\ref{lst:termination} expresses this theorem in Lean as
well-foundedness of the flipped one-step reduction relation.

\subsection{Translation to ideal congruence}

If $f\red{G}g$, then the definition gives
$f-g=cx^s p$ for some $p\in G$; hence $f-g\in\ideal{G}$.  The same conclusion
holds for a finite reduction sequence and for the equivalence closure of
reduction.

Conversely, assume Condition~\eqref{eq:unit-condition} and
$f-g\in\ideal{G}$.  Here \emph{span induction} means proving a property for
every element of $\ideal{G}$ by checking it first for each generator in $G$
and then showing that it is preserved by zero, addition, and polynomial
scalar multiplication.  For a polynomial $x$, let $Q(x)$ be the assertion
that
\[
  a+hx\eqvred{G}a
\]
for all polynomials $a$ and $h$.  If $p\in G$, then
$hp\redstar{G}0$; hence
\[
  (a+hp)-a=hp\redstar{G}0,
\]
and the difference-to-equivalence part of the polynomial translation lemma,
applied to $(a+hp)-a$, gives $a+hp\eqvred{G}a$.  In Lean, this implication is
\lean{MonomialOrder.eqvGen_of_sub_reflTransGen_zero}; the generator case is
packaged as \lean{MonomialOrder.eqvGen_add_mul_mem}.  Thus $Q$ holds for the
generators and is preserved by zero, addition, and polynomial scalar
multiplication.  Span induction therefore gives $Q(f-g)$, and the resulting
full equivalence is
\leandoc{Mathlib/RingTheory/MvPolynomial/PolynomialReductions}{MonomialOrder.eqvGen_iff_sub_mem_span}{MonomialOrder.eqvGen_iff_sub_mem_span}.
Taking
$a=g$ and $h=1$ yields
\[
  f=g+(f-g)\eqvred{G}g.
\]

\begin{lstlisting}[style=lean,caption={Reduction equivalence coincides with congruence modulo the generated ideal.},label={lst:translation}]
theorem eqvGen_iff_sub_mem_span
    (P : Set (MvPolynomial σ R))
    (hP₀ :
      ∀ p ∈ P,
        IsUnit (m.leadingCoeff p) ∨ p = 0)
    {f g : MvPolynomial σ R} :
    f ⟷*[m, P] g ↔ f - g ∈ Ideal.span P
\end{lstlisting}

\begin{theorem}[Reduction/ideal translation]\label{thm:translation}
Assume \eqref{eq:unit-condition}.  Then
\[
  f\eqvred{G}g
  \quad\Longleftrightarrow\quad
  f-g\in\ideal{G}.
\]
\end{theorem}

\section{Gr\"obner bases as confluent reduction systems}\label{sec:groebner-criterion}

Definition~\ref{def:groebner-conditions} distinguishes the self-relative
predicate
\leandoc{Mathlib/RingTheory/MvPolynomial/GroebnerBasisCriterion}{MonomialOrder.IsGroebner}{MonomialOrder.IsGroebner}
from the ideal-relative predicate
\leandoc{Mathlib/RingTheory/MvPolynomial/BuchbergerCriterion}{MonomialOrder.IsGroebnerBasis}{MonomialOrder.IsGroebnerBasis}.
The theorem
\leandoc{Mathlib/RingTheory/MvPolynomial/BuchbergerCriterion}{MonomialOrder.isGroebnerBasis_span_iff_isGroebner}{isGroebnerBasis_span_iff_isGroebner}
identifies them when
$I=\ideal{G}$.  Listing~\ref{lst:groebner-defs} displays the self-relative
Lean definition and its confluence characterization.

\begin{lstlisting}[style=lean,caption={The self-relative Gr\"obner condition and its confluence characterization.},label={lst:groebner-defs}]
def IsGroebner
    (G : Set (MvPolynomial σ R)) : Prop :=
  Ideal.span
      (m.leadingTerm ''
        (((Ideal.span G : Ideal (MvPolynomial σ R)) :
          Set (MvPolynomial σ R)))) =
    Ideal.span (m.leadingTerm '' G)

theorem IsGroebner.iff_confluent
    (G : Set (MvPolynomial σ R))
    (hG₀ :
      ∀ g ∈ G,
        IsUnit (m.leadingCoeff g) ∨ g = 0) :
    m.IsGroebner G ↔
      Relation.Confluent (m.ReducesToSet G)
\end{lstlisting}

\begin{theorem}[Reduction-theoretic characterizations of Gr\"obner bases]
\label{thm:groebner-characterizations}
Assume Condition~\eqref{eq:unit-condition}.  Then the following propositions
are equivalent:
\begin{enumerate}[label=(\alph*),leftmargin=2em]
  \item $\ideal{\LT(f)\mid f\in\ideal{G}}
        =\ideal{\LT(g)\mid g\in G}$;
  \item every $f\in\ideal{G}$ satisfies $f\redstar{G}0$;
  \item every nonzero element of $\ideal{G}$ is reducible modulo $G$;
  \item $\red{G}$ satisfies the Church--Rosser property;
  \item $\red{G}$ is locally confluent;
  \item $\red{G}$ is confluent;
  \item $\red{G}$ has unique normal forms; and
  \item each residue class modulo $\ideal{G}$ has a unique normal-form
        representative.
\end{enumerate}
\end{theorem}

\begin{table}[t]
\caption{Theorem~\ref{thm:groebner-characterizations} and its Lean declarations.}
\label{tab:groebner-characterizations-lean}
\centering
\scriptsize
\begin{tabularx}{\textwidth}{@{}>{\raggedright\arraybackslash}p{0.07\textwidth}>{\raggedright\arraybackslash}p{0.43\textwidth}>{\raggedright\arraybackslash}X@{}}
\toprule
Label & Mathematical condition & Lean declaration \doclinkmark \\
\midrule
(a) & Leading-term ideal criterion & \leandoc{Mathlib/RingTheory/MvPolynomial/GroebnerBasisCriterion}{MonomialOrder.IsGroebner}{MonomialOrder.IsGroebner} \\
(b) & Every element of $\ideal{G}$ reduces to zero & \leandoc{Mathlib/RingTheory/MvPolynomial/GroebnerBasisCriterion}{MonomialOrder.IsGroebner.iff_idealElements_reduceTo_zero}{iff_idealElements_reduceTo_zero} \\
(c) & Every nonzero element of $\ideal{G}$ is reducible & \leandoc{Mathlib/RingTheory/MvPolynomial/GroebnerBasisCriterion}{MonomialOrder.IsGroebner.iff_nonzero_idealElements_reducible}{iff_nonzero_idealElements_reducible} \\
(d) & Church--Rosser & \leandoc{Mathlib/RingTheory/MvPolynomial/GroebnerBasisCriterion}{MonomialOrder.IsGroebner.iff_churchRosser}{iff_churchRosser} \\
(e) & Local confluence & \leandoc{Mathlib/RingTheory/MvPolynomial/GroebnerBasisCriterion}{MonomialOrder.IsGroebner.iff_locallyConfluent}{iff_locallyConfluent} \\
(f) & Confluence & \leandoc{Mathlib/RingTheory/MvPolynomial/GroebnerBasisCriterion}{MonomialOrder.IsGroebner.iff_confluent}{iff_confluent} \\
(g) & Unique normal forms & \leandoc{Mathlib/RingTheory/MvPolynomial/GroebnerBasisCriterion}{MonomialOrder.IsGroebner.iff_uniqueNormalForms}{iff_uniqueNormalForms} \\
(h) & Unique normal-form representative in each residue class & \leandoc{Mathlib/RingTheory/MvPolynomial/GroebnerBasisCriterion}{MonomialOrder.IsGroebner.iff_unique_normalForm_representatives}{iff_unique_normalForm_representatives} \\
\bottomrule
\end{tabularx}
\medskip

\raggedright\scriptsize
The declarations corresponding to (b)--(h) are in the namespace
\lean{MonomialOrder.IsGroebner}.
\end{table}

The proof of Theorem~\ref{thm:groebner-characterizations} is organized as
explicit implications between the eight propositions.  Termination gives at
least one normal form reachable from every polynomial.  The equality in
(a) proves that a residue class modulo $\ideal{G}$ contains at most one
normal polynomial, and hence exactly one.  By Theorem~\ref{thm:translation},
equality of residue classes is the same as reduction equivalence, so uniqueness
of normal representatives yields Church--Rosser.
Theorem~\ref{thm:reduction-termination} and Newman's lemma identify local
confluence with confluence.  The Church--Rosser property implies that every
$f\in\ideal{G}$ reduces to zero because $f\eqvred{G}0$.

The converse from (c) to (a) requires more than an immediate
leading-term argument: a nonzero ideal element may initially be reducible at a
lower term.  The Lean proof uses well-founded induction on the support-colex
order.  If the reduction occurs below the leading term, the reduct has the same
leading term and a strictly smaller support-colex measure, so the induction
hypothesis applies.  Eventually one obtains a reduction at the leading term,
which shows that the leading term is divisible by the leading term of an
element of $G$ and proves the equality in (a).

\section{Buchberger's criterion}\label{sec:buchberger-criterion}

\begin{theorem}[Buchberger's criterion]
\label{thm:buchberger-criterion}
Assume Condition~\eqref{eq:unit-condition} and assume that $R$ has no zero
divisors.  Then
\[
  \ideal{\LT(f)\mid f\in\ideal{G}}
  =\ideal{\LT(g)\mid g\in G}
  \quad\Longleftrightarrow\quad
  \forall p,q\in G,\;\spoly(p,q)\redstar{G}0.
\]
\end{theorem}

The quantification in Theorem~\ref{thm:buchberger-criterion} includes zero and
diagonal pairs.  These cases add no substantive condition because
\[
  \spoly(0,q)=\spoly(p,0)=\spoly(p,p)=0,
\]
and the reflexive closure gives \(0\redstar{G}0\).  Moreover, every
polynomial that actually performs a one-step reduction is nonzero by
Definition~\ref{def:one-step-reduction}.

For the forward implication, every S-polynomial of two elements of $G$
belongs to $\ideal{G}$, so Theorem~\ref{thm:groebner-characterizations}
shows that it reduces to zero.  For the converse implication, assume that all
S-polynomials of pairs in $G$ reduce to zero.  The proof establishes local
confluence by analyzing a peak
\[
  f\red{G}f_1,
  \qquad
  f\red{G}f_2.
\]
Write the two reductions as $f_i=f-q_i$, where
$q_i=c_i x^{s_i} p_i$ cancels the term with exponent $t_i$.  If the two steps
cancel the same exponent of $f$, then there exist an exponent $v$ and a
coefficient $k\in R$ such that
\[
  f_2-f_1=q_1-q_2=kx^v\spoly(p_1,p_2).
\]
Finite reduction is preserved under multiplication by a monomial term, so the
hypothesis $\spoly(p_1,p_2)\redstar{G}0$ implies that
$f_2-f_1\redstar{G}0$.

If the two steps cancel different exponents, order the cancelled monomials,
say $x^{t_1}<_m x^{t_2}$.  The polynomial $q_1$ has leading exponent $t_1$,
so its coefficient at $t_2$ is zero.  Hence $q_1-q_2$ reduces in one step to
$q_1$, and $q_1=c_1x^{s_1} p_1$ reduces to zero.  The case
$x^{t_2}<_m x^{t_1}$ is symmetric.  Since
$f_2-f_1=q_1-q_2\redstar{G}0$, the difference-to-joinability part of the
polynomial translation lemma yields $f_1\joinred{G}f_2$.  Hence every one-step
peak is joinable, so $\red{G}$ is locally confluent.  The local-confluence equivalence from
Section~\ref{sec:groebner-criterion} then proves
\lean{MonomialOrder.IsGroebner}.

Listing~\ref{lst:buchberger-criterion} states both nontrivial assumptions:
every nonzero $g\in G$ has unit leading coefficient, and $R$ has no zero
divisors.  The no-zero-divisor assumption is used when computing leading terms
of polynomial products.

\begin{lstlisting}[style=lean,caption={Buchberger's criterion for the self-relative Gr\"obner predicate.},label={lst:buchberger-criterion}]
theorem IsGroebner.iff_sPolynomial_reflTransGen_zero
    [NoZeroDivisors R]
    (G : Set (MvPolynomial σ R))
    (hG₀ :
      ∀ g ∈ G,
        IsUnit (m.leadingCoeff g) ∨ g = 0) :
    m.IsGroebner G ↔
      ∀ g₁ ∈ G, ∀ g₂ ∈ G,
        m.sPolynomial g₁ g₂ ⟶*[m, G] 0
\end{lstlisting}

The file also contains an ideal-relative form.  If
$\operatorname{Ideal.span}(G)=I$, then
\lean{MonomialOrder.IsGroebnerBasis} is equivalent to the statement that every
S-polynomial of two elements of $G$ reduces to zero modulo $G$.

\section{An abstract Buchberger procedure}\label{sec:buchberger-procedure}

The Buchberger procedure is formalized as a relation rather than as a recursive function.

\begin{definition}[Buchberger step and the S-polynomial condition]
\label{def:buchberger-step}
A Buchberger step from $G$ to $H$ consists of $p,q\in G$ and a polynomial $h$
such that
\begin{enumerate}[label=(\roman*),leftmargin=2em]
  \item $h$ is a normal form of $\spoly(p,q)$ modulo $G$;
  \item $h\neq0$ and $\LC(h)$ is a unit; and
  \item $H=G\cup\{h\}$.
\end{enumerate}
We say that $G$ satisfies the \emph{S-polynomial condition} if
$\spoly(p,q)\redstar{G}0$ for all $p,q\in G$.  The corresponding Lean
declarations are
\leandoc{Mathlib/RingTheory/MvPolynomial/BuchbergerAlgorithm}{MonomialOrder.BuchbergerStep}{BuchbergerStep}
and
\leandoc{Mathlib/RingTheory/MvPolynomial/BuchbergerAlgorithm}{MonomialOrder.SPolynomialCondition}{SPolynomialCondition}.
The step relation does not prescribe
which pair $(p,q)$ or which normal form $h$ must be chosen.
\end{definition}

Although Definition~\ref{def:buchberger-step} does not explicitly require
$p$ and $q$ to be nonzero and distinct, every genuine Buchberger step has
witnesses with those properties.  If $p=0$, $q=0$, or $p=q$, then
$\spoly(p,q)=0$.  The zero polynomial has no outgoing reduction, so its only
reachable normal form is zero; this contradicts the requirement $h\neq0$ in a
Buchberger step.

\begin{lstlisting}[float=htbp,style=lean,caption={The nondeterministic Buchberger step relation.},label={lst:buchberger-step}]
def BuchbergerStep
    (G H : Set (MvPolynomial σ R)) : Prop :=
  ∃ p ∈ G, ∃ q ∈ G, ∃ h : MvPolynomial σ R,
    Relation.IsNormalFormOf
        (m.ReducesToSet G)
        (m.sPolynomial p q) h ∧
      h ≠ 0 ∧
        IsUnit (m.leadingCoeff h) ∧
          H = insert h G

def SPolynomialCondition
    (G : Set (MvPolynomial σ R)) : Prop :=
  ∀ p ∈ G, ∀ q ∈ G,
    m.sPolynomial p q ⟶*[m, G] 0
\end{lstlisting}

\begin{proposition}[Ideal preservation and strict leading-term growth]
\label{prop:buchberger-growth}
Let $G,H\subseteq R[X_\sigma]$.
\begin{enumerate}[label=(\roman*),leftmargin=2em]
  \item If there is a Buchberger step from $G$ to $H$, then
  \[
    \ideal{H}=\ideal{G}.
  \]
  \item Assume that every nonzero element of $G$ has unit leading
  coefficient.  If there is a Buchberger step from $G$ to $H$, then
  \[
    \ideal{\LT(g)\mid g\in G}
    \subsetneq
    \ideal{\LT(h')\mid h'\in H}.
  \]
\end{enumerate}
\end{proposition}

For part~(i), the adjoined polynomial $h$ is a reduction result of an
S-polynomial of elements of $G$, and therefore $h\in\ideal{G}$.  For
part~(ii), suppose that $\LT(h)$ belonged to the ideal generated by the
leading terms of $G$.  Under the unit-leading-coefficient hypothesis, this
membership produces a one-step reduction of $h$, contradicting the fact that
$h$ is a nonzero normal form modulo $G$.  The subtype
\leandoc{Mathlib/RingTheory/MvPolynomial/BuchbergerAlgorithm}{MonomialOrder.BuchbergerState}{BuchbergerState}
packages precisely this invariant: every nonzero
polynomial in the current set has unit leading coefficient.

\begin{lstlisting}[style=lean,caption={Noetherian termination of the Buchberger step relation.},label={lst:buchberger-termination}]
def BuchbergerState : Type _ :=
  { G : Set (MvPolynomial σ R) //
      ∀ g ∈ G,
        IsUnit (m.leadingCoeff g) ∨ g = 0 }

def BuchbergerState.Step
    (G H : BuchbergerState m) : Prop :=
  m.BuchbergerStep G.1 H.1

theorem BuchbergerState.wellFounded_flip
    [IsNoetherianRing (MvPolynomial σ R)] :
    WellFounded
      (flip (BuchbergerState.Step (R := R) (m := m)))
\end{lstlisting}

The termination measure sends a Buchberger state $G$ to the ideal
$\ideal{\LT(g)\mid g\in G}$.  Because the state contains the
unit-leading-coefficient invariant, every step between Buchberger states
strictly increases this ideal by Proposition~\ref{prop:buchberger-growth}(ii).
Since a Noetherian ring admits no infinite strictly ascending chain of ideals,
there is no infinite forward chain of Buchberger-state steps.  In Lean, this
is encoded as well-foundedness of the flipped step relation, as shown in
Listing~\ref{lst:buchberger-termination}.

\begin{theorem}[Existence and correctness of a terminal Buchberger state]
\label{thm:terminal-correctness}
Let $G\subseteq R[X_\sigma]$ and let $I=\ideal{G}$.  Assume that $R$ has no
zero divisors, $R[X_\sigma]$ is Noetherian, and every nonzero polynomial in
$R[X_\sigma]$ has unit leading coefficient.  Then no infinite sequence of
Buchberger steps starts from $G$.  Moreover, there exist a Buchberger state
$H$ and a finite sequence of Buchberger steps from $G$ to $H$ such that $H$
has no outgoing Buchberger step and the underlying polynomial set of $H$ is a
Gr\"obner basis of $I$.
\end{theorem}

The Lean theorem also takes the corresponding initial-state invariant as an
explicit hypothesis.  The global coefficient hypothesis in
Theorem~\ref{thm:terminal-correctness} implies this invariant.

If $R=K$ is a field, the no-zero-divisor and unit-leading-coefficient
hypotheses in Theorem~\ref{thm:terminal-correctness} are automatic.  If $K$ is
Noetherian and the index type $\sigma$ is finite, Mathlib's form of
Hilbert's basis theorem proves that \texttt{MvPolynomial }\(\sigma\)\texttt{ K} is Noetherian.

\begin{corollary}[Finite terminal Gr\"obner basis over a field]
\label{cor:finite-terminal}
Let $K$ be a Noetherian field, let $\sigma$ be finite, and let
$G\subseteq K[X_\sigma]$ be finite.  If $I=\ideal{G}$, then there exist a
Buchberger state $H$ whose underlying polynomial set is finite and a finite sequence of Buchberger steps from $G$
to $H$ such that $H$ has no outgoing step and its underlying polynomial set is
a Gr\"obner basis of $I$.
\end{corollary}

The finite-output result follows from two preservation theorems in the Lean
development: one Buchberger step preserves finiteness, and every finite chain
of steps preserves finiteness.  These results are
\leandoc{Mathlib/RingTheory/MvPolynomial/BuchbergerAlgorithm}{MonomialOrder.BuchbergerStep.finite}{MonomialOrder.BuchbergerStep.finite}
and
\leandoc{Mathlib/RingTheory/MvPolynomial/BuchbergerAlgorithm}{MonomialOrder.ReflTransGen_BuchbergerStep.finite}{ReflTransGen_BuchbergerStep.finite},
respectively.

\begin{remark}[What is and is not executable]\label{rem:not-executable}
The finiteness results above do not make the construction executable: the
development does not define a \lean{Finset}-based state, a finite queue of
critical pairs, a pair-selection function, or a normal-form function.
\end{remark}

\section{Formalization design and reproducibility}\label{sec:design}

\subsection{Relation versus division function}

Mathlib already provides polynomial division with respect to a monomial
order.  The relation $\red{G}$ has a different role: it specifies which
single cancellations are mathematically valid, independently of how an
algorithm chooses them.  The current development proves bridge lemmas showing
that a finite relational reduction sequence yields a finite linear-combination
representation of the input, together with the same product-degree bound that
appears in \lean{MonomialOrder.div_set}.  It does not yet prove that the remainder
returned by Mathlib's division procedure is reachable by $\redstar{G}$ and is
normal modulo $G$; establishing that direct connection remains future work.
The relational specification can also be used to compare leading-term,
arbitrary-term, matrix-based, and signature-restricted reduction procedures.

\subsection{Existing and new infrastructure}

The modified \lean{Relation.lean} file contains both pre-existing Mathlib
code and new declarations.  The new declarations include \lean{Diamond},
\lean{Confluent}, \lean{ChurchRosser}, the closure bridges
\lean{Relation.ReflTransGen.to_eqvGen} and \lean{Relation.Join.to_eqvGen}, and
the proofs relating confluence and Church--Rosser.  The file \lean{NormalForm.lean} is separate and defines the
normal-form API used by the polynomial files.  Accordingly, the paper does not
count the complete length of \lean{Relation.lean} as newly formalized code.

\subsection{General coefficient rings}

The coefficient assumptions are layered according to the result being proved.
\begin{itemize}[leftmargin=2em]
  \item Over an arbitrary commutative ring, the development defines reduction,
    characterizes reducibility and normality, proves termination and degree
    bounds, and extracts finite linear-combination certificates such as
    \leandoc{Mathlib/RingTheory/MvPolynomial/PolynomialReductions}{MonomialOrder.exists_linearCombination_of_reflTransGen_with_degree_bound}{exists_linearCombination_of_reflTransGen_with_degree_bound}.
  \item To prove that an arbitrary polynomial multiple of a reducer reduces to
    zero, it is enough for the reducer's leading coefficient to be regular,
    that is, a non-zero-divisor.
  \item The translation lemma, the equivalence between reduction congruence and
    ideal congruence, and the principal confluence/Gr\"obner equivalences use
    the stronger unit-leading-coefficient hypothesis
    \eqref{eq:unit-condition}.
\end{itemize}
Over a coefficient field, Condition~\eqref{eq:unit-condition} is automatic.

\subsection{Reproducibility of the formal artifact}
\label{sec:reproducibility}

The six reviewed Lean files contain no occurrences of \lean{sorry} or
\lean{admit}.  Two immutable snapshots identify the reviewed development:
\begin{itemize}[leftmargin=2em]
  \item the canonical Lean source is commit
    \href{https://github.com/Sanghyeok0/mathlib4/tree/4868fec773dcba035d283904f8b118511cc83d73}
    {\lean{4868fec773dcba035d283904f8b118511cc83d73}}
    of the author's \lean{mathlib4} fork;
  \item the companion artifact is commit
    \href{https://github.com/Sanghyeok0/Buchberger/commit/3e7bbc935df0f545f4d56ed4a457fc23510e30e7}
    {\lean{3e7bbc935df0f545f4d56ed4a457fc23510e30e7}}
    of the \lean{Buchberger} repository.
\end{itemize}
The artifact pins the source commit and the Lean
\lean{v4.33.0-rc1} toolchain.  A clean reproduction uses
\begin{verbatim}
git clone https://github.com/Sanghyeok0/Buchberger.git
cd Buchberger
git checkout 3e7bbc935df0f545f4d56ed4a457fc23510e30e7
lake build
\end{verbatim}
The root module built by this command contains no copied formalization
code: it imports exactly the six modules shown in
Figure~\ref{fig:dependencies} from the pinned fork.  For a direct check of the
final module in the canonical source repository, run
\begin{verbatim}
git clone https://github.com/Sanghyeok0/mathlib4.git
cd mathlib4
git checkout 4868fec773dcba035d283904f8b118511cc83d73
lake exe cache get
lake env lean Mathlib/RingTheory/MvPolynomial/BuchbergerAlgorithm.lean
\end{verbatim}
The artifact homepage
\url{https://sanghyeok0.github.io/Buchberger/} provides reproduction
instructions.  Generated API documentation is available at
\url{https://sanghyeok0.github.io/Buchberger/docs/}; each documented
declaration links to its definition in the fixed \lean{mathlib4} source.

\section{Related work}\label{sec:related-work}

\citet{GuoEtAl2026} formalized polynomial division with remainder,
Buchberger's criterion, reduced Gr\"obner bases, and results connecting
finite-variable and infinite-variable polynomial rings in Lean 4.  Therefore,
this article does not claim to be the first Lean formalization of Gr\"obner
basis theory.  The ideal-relative predicate \lean{IsGroebnerBasis} is a
standard definition also used in that development; this article does not
claim the definition itself as a new contribution.  In comparison, the
present development focuses on the explicit arbitrary-term reduction relation; the
equivalence between the Gr\"obner condition and Church--Rosser, local
confluence, confluence, reduction to zero, and uniqueness of normal forms;
and the well-founded nondeterministic Buchberger step relation.

\citet{ShenEtAl2026} developed tactics that send polynomial computations to
SageMath or SymPy and verify the returned certificates in Lean using a
computable polynomial representation.  Their system is designed for practical
automation on concrete inputs.  The formalization in this article does not call
an external computer algebra system: it proves the reduction theory and theorems about the abstract Buchberger
procedure inside Lean.  Conversely, it does not yet compute Gr\"obner bases on
concrete inputs.

In Isabelle/HOL, \citet{MaletzkyImmler2018} formalized Gr\"obner bases of
modules and executable implementations of Buchberger's and Faug\`ere's $F_4$
algorithms.  \citet{Maletzky2021} then formalized a generic executable family
of signature-based algorithms.  Those Isabelle developments include concrete
functions that can be evaluated, whereas \lean{BuchbergerStep} in this article
is a proposition-valued relation.  This difference motivates the future
separation between the relation \lean{BuchbergerStep} and an efficient
ordered representation for executable Lean code.

\section{Limitations and future work}\label{sec:future-work}

The proved theorems provide a relational specification for polynomial reduction and the abstract Buchberger procedure.  The following results and implementations are not yet part of the development.
\begin{itemize}[leftmargin=2em]
  \item Define an executable Buchberger procedure using \lean{Finset}-based
    states, a finite queue of critical pairs, a computable pair-selection
    strategy, and a verified normal-form function.
  \item Connect Mathlib's existing division procedure to the relational normal
    form specification.
  \item Separate specification-level polynomials from a computable ordered
    representation suitable for efficient execution and reflection.
  \item Formalize interreduction and reduced Gr\"obner bases in the
    reduction-theoretic interface.
  \item Extend the framework to matrix-based $F_4$ reduction and to
    signature-based algorithms, including $F_5$-style criteria.
\end{itemize}

\section{Conclusion}\label{sec:conclusion}

The formalization proves a specific chain of results.  Arbitrary-term
polynomial reduction is terminating over an arbitrary commutative coefficient
ring.  Under the unit-leading-coefficient hypothesis, reduction equivalence
modulo $G$ agrees with congruence modulo $\ideal{G}$, and the Gr\"obner
condition is equivalent to reduction-to-zero and the confluence properties.
Assuming additionally that the coefficient ring has no zero divisors,
Buchberger's S-polynomial condition is equivalent to the Gr\"obner condition.
On the state space carrying the unit-leading-coefficient invariant,
Noetherianity prevents infinite sequences of Buchberger-state steps.  One step
and finite chains preserve finiteness, so finite input yields a finite terminal
Gr\"obner basis in the finite-variable field specialization.  These theorems
provide a verified relational specification of the abstract Buchberger
procedure.  A practical certified solver additionally requires finite
executable data structures, a verified normal-form function, and a computable
strategy for selecting polynomial pairs and processing their S-polynomials.

\section*{Acknowledgments}
The author thanks the Lean and Mathlib communities for the existing library,
discussion, and review on which this work builds.

\section*{Funding}
This work was supported by the National Research Foundation of Korea (NRF)
grant funded by the Korea government (MSIT) (No. RS-2025-00520280).

\section*{Data and software availability}
No external dataset was generated or analyzed.  The canonical Lean source is
publicly available in the author's \lean{mathlib4} fork.  The reproducible
companion artifact, CI configuration, and generated documentation are
available in the \lean{Buchberger} repository and through the artifact homepage
\url{https://sanghyeok0.github.io/Buchberger/}.  The exact source and artifact
snapshots, Lean toolchain, and reproduction commands are recorded in
Section~\ref{sec:reproducibility}.


\small
\setlength{\bibsep}{2pt}
\bibliographystyle{plainnat}
\bibliography{references}

\end{document}
